\documentclass[12pt]{article}
\usepackage[utf8]{inputenc}
\usepackage[T1]{fontenc}
\usepackage[margin=1in]{geometry}
\usepackage{setspace}
\usepackage{xcolor}
\usepackage[colorlinks=true, urlcolor=blue!60!black, linkcolor=blue!60!black]{hyperref}
\setlength{\parindent}{2em}
\setlength{\parskip}{1.25em}
\renewcommand{\baselinestretch}{1.0}
\begin{document}
\begin{center}
{\LARGE \textit{Buckley v.\ Valeo}, 424 U.S.\ 1 (1976)\par}
\vspace{0.8em}
{\large Case Brief --- Four-Page Edition\par}
\end{center}
\vspace{1em}

\section{Memory Jogger}

The foundation of modern campaign-finance law: contribution limits, disclosure, and voluntary presidential public financing upheld; limits on campaign and independent expenditures struck down under the First Amendment; congressional appointment of Federal Election Commission members struck down under the Appointments Clause (\href{https://supreme.justia.com/cases/federal/us/424/1/}{Justia}).

\section{Detailed Case Facts}

After Watergate, Congress enacted the 1974 amendments to the Federal Election Campaign Act (FECA). The Act (1) capped contributions at \$1,000 per candidate per election (\$5,000 for political committees; \$25,000 per year in total), capped independent expenditures ``relative to a clearly identified candidate'' at \$1,000 a year, capped candidates' personal and family spending, and imposed overall campaign ceilings; (2) required recordkeeping for contributions over \$10 and public disclosure of contributors over \$100; (3) created public financing of presidential campaigns, funded by a tax check-off and divided among major, minor, and new parties, conditioned on spending caps; and (4) created a Federal Election Commission (FEC) with rulemaking and enforcement powers, four of whose six voting members were appointed by congressional leaders. An ideologically diverse coalition---Senator James L.\ Buckley, former Senator Eugene McCarthy, the New York Civil Liberties Union, the Libertarian Party, and others---sued Senate Secretary Francis Valeo, the FEC, and other officials.

\section{Procedural History}

The en banc D.C.\ Circuit upheld nearly all of FECA, relying on \textit{United States v.\ O'Brien} to treat money as regulable conduct; a three-judge district court upheld public financing. On direct appeal, the Supreme Court issued a lengthy \textit{per curiam} opinion affirming in part and reversing in part.

\section{Judicial Votes}

Eight Justices participated (Stevens took no part; Douglas had retired). Brennan, Stewart, and Powell joined the \textit{per curiam} in full; the others joined parts.
\begin{itemize}
\item \textbf{Contribution limits upheld:} 6--2 (Burger and Blackmun dissenting).
\item \textbf{Independent and overall expenditure limits struck:} 7--1 (White dissenting).
\item \textbf{Personal-funds limit struck:} 6--2 (Marshall and White dissenting).
\item \textbf{Disclosure upheld:} 7--1 (Burger dissenting in part).
\item \textbf{Public financing upheld:} 6--2 on general-election financing (Burger and Rehnquist dissenting).
\item \textbf{FEC appointments struck:} 7--1 (White dissenting).
\end{itemize}

\section{Holding}

Political spending is speech, not mere conduct. Contribution limits impose only a marginal restriction and survive if closely drawn to the important interest in preventing \textit{quid pro quo} corruption or its appearance. Expenditure limits directly restrain the quantity of political speech, trigger exacting scrutiny, and cannot be justified by anti-corruption or equalization interests. Independent-expenditure and reporting rules reach only express advocacy (``vote for,'' ``elect''). Disclosure is valid, with as-applied exemptions for groups showing a ``reasonable probability'' of harassment. Voluntary public financing is valid. Anyone exercising significant authority under federal law is an Officer of the United States who must be appointed under Article~II, so the FEC could not exercise enforcement powers.

\section{Analysis of Opinions}

\textbf{Per Curiam.} Because ``virtually every means of communicating ideas in today's mass society requires the expenditure of money,'' spending limits restrict speech. A contribution, however, is a ``general expression of support'' whose communicative value does not grow with its size, and large contributions risk corruption; courts lack a ``scalpel'' to second-guess the dollar amount. Independent expenditures pose a ``substantially diminished potential for abuse,'' and self-funding reduces dependence on donors. Disclosure serves voter information, corruption deterrence, and enforcement. Public financing enlarges speech. Congress cannot aggrandize itself by appointing officers who execute federal law.

\textbf{Burger (in part).} ``Contributions and expenditures are two sides of the same First Amendment coin''; he would have struck contribution limits, small-contribution disclosure, and public financing.

\textbf{White (in part).} ``Money is not speech''; Congress deserved deference on expenditure ceilings and on FEC appointments.

\textbf{Marshall (in part).} Limiting outside contributions while letting wealthy candidates spend their fortunes gives the rich an asymmetric advantage.

\textbf{Blackmun and Rehnquist (in part).} Blackmun saw no principled line between contributions and expenditures. Rehnquist argued the public-financing formula entrenched the two major parties.

\section{Examples: Future Applications}

\begin{enumerate}
\item \textbf{Same-Side: State Contribution Cap.} A \$1,075 per-election cap is constitutional; courts do not act as a ``scalpel'' on the exact amount (\textit{Nixon v.\ Shrink Missouri}).
\item \textbf{Same-Side: Independent PAC Ads.} A \$5,000 cap on uncoordinated PAC ads is unconstitutional (\textit{FEC v.\ NCPAC}).
\item \textbf{Opposite-Side: Aggregate Limits.} Although \textit{Buckley} upheld a \$25,000 aggregate cap, its \textit{quid pro quo} framework now invalidates aggregate limits (\textit{McCutcheon v.\ FEC}).
\item \textbf{Opposite-Side: Harassed Minor Party.} A party with documented surveillance, firings, and threats meets the ``reasonable probability'' standard and is exempt from donor disclosure.
\item \textbf{Fence-Sitter: Trigger Matching Funds.} Extra public funds triggered by an opponent's spending add speech yet penalize expenditures; \textit{Arizona Free Enterprise Club v.\ Bennett} (2011) struck them down.
\end{enumerate}

\section{Critique}

\textbf{An unstable dichotomy.} Critics across the spectrum say limiting the supply of money while freeing demand channeled billions into super PACs and ``dark money,'' and the footnote-52 ``magic words'' test was easily evaded by issue ads.

\textbf{From both flanks.} Progressives argue \textit{Buckley} equated money with speech and disabled Congress from addressing wealth's distortion of politics. Originalists such as Justice Thomas argue the First Amendment draws no line between contributions and expenditures, so reduced scrutiny lets incumbents write self-serving caps, while praising Part IV's Appointments Clause formalism.

\section{Key Quotations}

\begin{quote}
``But the concept that government may restrict the speech of some elements of our society in order to enhance the relative voice of others is wholly foreign to the First Amendment, which was designed `to secure the widest possible dissemination of information from diverse and antagonistic sources.'\,''
\end{quote}
\noindent --- Per Curiam, 424 U.S.\ at 48--49.

\begin{quote}
``We think its fair import is that any appointee exercising significant authority pursuant to the laws of the United States is an `Officer of the United States,' and must, therefore, be appointed in the manner prescribed by \S~2, cl.~2, of that Article.''
\end{quote}
\noindent --- Per Curiam, 424 U.S.\ at 126.

\begin{quote}
``By contrast with a limitation upon expenditures for political expression, a limitation upon the amount that any one person or group may contribute to a candidate or political committee entails only a marginal restriction upon the contributor's ability to engage in free communication.''
\end{quote}
\noindent --- Per Curiam, 424 U.S.\ at 20--21.

\end{document}
